\section{Success--Completion-Time Pareto}
\label{appx:success_pareto}

Figure~\ref{fig:s1_pareto_success} complements the recovery-focused Pareto in the main text with the task-completion view. Each method is plotted by its held-out success rate ($x$-axis, higher is better) and its mean completion time $t_{\mathrm{succ}}$ ($y$-axis, lower is better); error bars are sample standard deviations over the three training seeds.

EA-RG lies close to the high-success frontier with a low completion time ($0.985 \pm 0.011$ success, $t_{\mathrm{succ}} = 46.1 \pm 0.4$ steps). The wider single-graph baseline and HAPPO reach slightly higher success rates ($0.997$ and $1.000$) but complete tasks more slowly ($57.6 \pm 18.0$ and $49.9 \pm 3.0$ steps). This figure supports the task-level efficiency claim from a second angle and is not intended to repeat the recovery-latency conclusion of RQ2.

\begin{figure}[t]
    \centering
    \includegraphics[width=0.9\linewidth]{figures/fig_pareto_success.png}
    \caption{Success--completion-time Pareto (n=3 training seeds; error bars are sample SD). 
    EA-RG is close to the high-success frontier with the lowest completion time among the 
    first-tier methods.}
    \label{fig:s1_pareto_success}
\end{figure}
